\centerline{\bf \Huge Домашнее задание}
\centerline{\bf \Huge  Начало графов}


\begin{flushright}
        \scriptsize{Те, кто вместе идут к победе $-$ и есть друзья. \\Наруто}
\end{flushright}


\problem{1} В ЭлМШ любой ученик знаком с любым другим через какое-то количество рукопожатий. Известно, что в ЭлМШ обучаются 40 учеников, и на всех приходится 39 дружб. Может ли существовать треугольник дружбы (Федор с Игорем, Игорь с Элианой, Элиана с Федором)?
  
\problem{2} Как известно, в Республике Калмыкия 13 районов, а в Республике Дагестан их целых 41! Представим, что в Калмыкии 11 дорог, каждая из которых соединяет два района внутри Калмыкии, а в Дагестане таких дорог 39. Однажды Начальник и НеНачальник решили вложить маленькую часть бюджета ЭлМШ в проложение двух дорог между двумя республиками(от двух районов одной к двум районам другой). Может ли теперь Начальник из Калмыкии приехать в гости к НеНачальнику в Дагестан?

\problem{3} В Элисте десять микрорайонов, и каждая дорога в городе соединяет какие-то два микрорайона. Может ли из шестого микрорайона выходить шесть дорог, из седьмого - семь, а из восьмого - восемь?